\documentclass[11pt, a4paper]{article}

\usepackage[T1]{fontenc}
\usepackage[utf8]{inputenc}
\usepackage{tgpagella}
\usepackage[spanish, es-noshorthands]{babel}
\usepackage{amsmath, amssymb, bm}
\usepackage{tabularx}
\usepackage{booktabs}
\usepackage[most]{tcolorbox}
\usepackage[margin=2.5cm]{geometry}
\usepackage{ragged2e}
\usepackage{fancyhdr}

\pagestyle{fancy}
\fancyhf{}
\setlength{\headheight}{16pt}
\lhead{\textbf{UCB Sede Tarija} | Ing. Mecatrónica}
\chead{}
\rhead{IMT-121 Circuitos Electrónicos I | BOM y Control de Aprobación}
\lfoot{Prof: M.Sc. Ing. Bernardo Quiroga Turdera}
\rfoot{Página \thepage}
\renewcommand{\headrulewidth}{0.4pt}

\definecolor{ucbblue}{RGB}{0, 51, 102}
\definecolor{warnred}{RGB}{200, 0, 0}

\begin{document}

\begin{tcolorbox}[colback=warnred!5!white, colframe=warnred, title=\textbf{BOM PROVISIONAL Y GATE DE AUTORIZACIÓN (USO DOCENTE)}]
    \centering \textbf{Preparación de Laboratorio Integrado: RLC Serie} \\
    \textbf{¡IMPORTANTE: NO ORDENAR COMPRAS HASTA VALIDAR ESTA HOJA!}
\end{tcolorbox}

\section*{1. Modelo Operativo y Rotación Asumida}
El diseño de la experiencia depende estrictamente del inventario funcional real.
\begin{itemize}
    \item \textbf{Recursos Tentativos:} 3 Generadores de Funciones y 3 Osciloscopios.
    \item \textbf{Esquema de Rotación:} Cohorte de 9 grupos distribuidos en \textbf{3 Tandas (3 grupos por turno)}. Un cuarto osciloscopio no está confirmado ni asegurado en este planeamiento.
\end{itemize}

\section*{2. Lista de Materiales y Equipos (BOM Candidata)}

\renewcommand{\arraystretch}{1.5}
\noindent
\begin{tabularx}{\textwidth}{>{\bfseries\RaggedRight}p{3cm} >{\RaggedRight\arraybackslash}X >{\centering\arraybackslash}p{2.5cm} >{\centering\arraybackslash}p{2.5cm}}
\toprule
\textbf{Ítem y Función} & \textbf{Especificación Técnica a Comprobar} & \textbf{Cantidad} & \textbf{Acción Pdt.} \\ 
\midrule
Generador de Señales & Salida efectiva, impedancia interna ($R_{gen}$), rise time, amplitud nominal y masa. & 3 pares con el Osciloscopio & Probar 1 a 1 \\ \addlinespace
Osciloscopio & Base de tiempo, tasa de muestras, sondas de medición, modo de disparo y \textbf{aislamiento}. & 3 equipos. & Probar 1 a 1 \\ \addlinespace
Inductor (Energía Magnética) & \textbf{Candidato: $10\,\text{mH}$}. Verificar resistencia DC ($R_{DCR}$), tolerancia y corriente de saturación. & Según \# Grupos & NO ORDENAR \\ \addlinespace
Capacitor (Energía Eléctrica) & \textbf{Candidato: $1\,\mu\text{F}$ No Polarizado}. Verificar límite dieléctrico/tensión, tolerancia. & Según \# Grupos & NO ORDENAR \\ \addlinespace
Resistores & Para modelar $R_{total}$ sub/crítico/sobreamortiguado tras deducir $R_{gen} + R_{DCR}$. & Pendiente & NO ORDENAR \\ 
\bottomrule
\end{tabularx}

\vspace{0.6cm}
\begin{tcolorbox}[colback=gray!5, colframe=black, title=\textbf{Gate de Autorización Docente y Seguridad (Hito de Validación)}]
La Práctica Integrada NO ESTÁ AUTORIZADA hasta que el Docente Responsable firme o tilde todas las comprobaciones ineludibles:
\begin{itemize}
    \item[$\square$] Probados y operativos 3 Generadores y 3 Osciloscopios con sus accesorios.
    \item[$\square$] \textbf{Crítico:} Revisadas las referencias de tierra. Ninguna pinza de masa será conectada a nodos activos (potenciales flotantes) que cortocircuiten el generador.
    \item[$\square$] Deducidos mediante óhmetro los valores de $R_{DCR}$ del inductor real.
    \item[$\square$] Calculados los regímenes de disipación de potencia de los resistores a adquirir para evitar quemaduras o humo.
    \item[$\square$] Ejecutada una prueba "Piloto" por parte del docente con el montaje final (generar, capturar, comparar $\omega_0$ y medir el tiempo logístico real de la prueba).
\end{itemize}
\end{tcolorbox}

\end{document}
